\chapter{Symmetry in Quantum Physics}

\section{Symmetry Groups in Quantum Physics}

\textbf{Symmetry transformations} are operations carried out on the devices which prepare the states and the apparatuses which measure the observables of a quantum system having the property of not changing the results of experiments. For instance, space translations are symmetry transformations since the uniform translation of all preparations devices and measuring apparatuses does not change the results of experiments due to the homogeneity of space. Likewise, space rotations are symmetry transformations since the uniform rotation of all preparations devices and measuring apparatuses does not change the results of experiments due to the isotropy of space. In relativity, Lorentz boosts are symmetry transformations since as the result of experiments are the same in a given inertial frame and any inertial frame moving at uniform speed with respect to the first because the isotropy of spacetime.

In general, symmetry transformations gather together in families with special properties. Let \(G\) be one such family. Carrying out successively two symmetry transformations \(g\), \(g'\) of \(G\) yields a third symmetry transformation \(g' \cdot g\) of \(G\) called the \textbf{product} of \(g\), \(g'\). The effect of a symmetry transformation \(g\) of \(G\) can be undone by another symmetry transformation \(g^{-1}\) of \(G\) called the \textbf{inverse} of \(g\). There exists a distinguished symmetry transformation \(e\) of \(G\) that has no effect called \textbf{neutral}. The multiplication and inversion operations and neutral element of \(G\) have the following properties:
\begin{itemize}
    \item \textbf{Associativity}: for any three symmetry transformations \(g\), \(g'\), \(g''\) of \(G\) we have
          \begin{equation}
              (g'' \cdot g') \cdot g = g'' \cdot (g' \cdot g).
              \label{eq:group_associativity_definition}
          \end{equation}
    \item \textbf{Neutrality}: there exists a symmetry transformation \(e\) of \(G\) such that for any symmetry transformation \(g\) of \(G\) we have
          \begin{equation}
              e \cdot g = g \cdot e = g.
              \label{eq:group_neutrality_definition}
          \end{equation}
    \item \textbf{Invertibility}: for any symmetry transformation \(g\) of \(G\) there exists a symmetry transformation \(g^{-1}\) of \(G\) such that
          \begin{equation}
              g^{-1} \cdot g = g \cdot g^{-1} = e.
              \label{eq:group_invertibility_definition}
          \end{equation}
\end{itemize}
These relations characterize \(G\) as an algebraic structure called a \textbf{group}. Since \(G\) is constituted by symmetry transformations, \(G\) is called more specifically a \textbf{symmetry group}. A group \(G\) is called \textbf{finite} if it is constituted by a finite number of elements, else it is called \textbf{infinite}. A group \(G\) is called \textbf{Abelian} if multiplication turns out to be commutative, that is
\begin{equation}
    g \cdot g' = g' \cdot g, \quad \forall g, g' \in G.
    \label{eq:group_abelian_definition}
\end{equation}
Else, it is called non Abelian. The groups encountered in quantum physics are bot finite and infinite and Abelian and non Abelian.

\begin{example}[Space Translation Symmetry]
    [...]\TODO{Complete the example}
\end{example}

\begin{example}[Space Rotation Symmetry]
    [...]\TODO{Complete the example}
\end{example}

\subsection{Symmetry Groups and Their Operator Realization}

In quantum theory, states and observables of a system are represented respectively by normalized vectors defined up to a multiplicative phase in a Hilbert space \(\mathcal{H}\) and selfadjoint operators in the operator algebra \(L(\mathcal{H})\). The invariance of experimental results under the action of a symmetry group reduces to the invariance of expectation values: if \(\Psi\) is a state, \(O\) is an observable and \(g\) is a symmetry transformation of \(G\), \(^g \Psi\) is the transformed state of \(\Psi\) and \(^g O\) is the transformed observable of \(O\), one must have
\begin{equation}
    \bra{^g \Psi} ^g O \ket{^g \Psi} = \bra{\Psi} O \ket{\Psi}.
    \label{eq:group_invariance_expectation_value_definition}
\end{equation}

Let us now analyze the implications of the invariance of expectation values under the action of a symmetry group \(G\) with a simple example.

\begin{example}[``yes–no'' observables]
    As is well–known, there is a class of elementary observables, the “yes–no” observables, defined by taking the value 1, 0 according to whether the system is in a certain state \(\Phi\) or not. The associated operator is given by
    \[
        O_{\Phi} \Psi = \Phi \bra{\Phi} \ket{\Psi}.
    \]
    It is natural to assume that under the symmetry transformation \(g\) of \(G\) the observable \(O_{\Phi}\) transforms accordingly with the state \(\Phi\) as follows:
    \[
        ^g O_{\Phi} = O_{^g \Phi}.
    \]
    The expectation value of \(O_{\Phi}\) in the state \(\Psi^{\prime}\) is given by the transition probability from \(\Psi^{\prime}\) to \(\Phi\):
    \[
        \bra{\Psi^{\prime}} O_{\Phi} \ket{\Psi^{\prime}} = |\bra{\Phi} \ket{\Psi^{\prime}}|^2,
    \]
    such that the \textbf{invariance of expectation values} under the action of the symmetry group \(G\) implies that
    \[
        |\bra{^g \Phi} \ket{^g \Psi^{\prime}}|^2 = |\bra{\Phi} \ket{\Psi^{\prime}}|^2.
    \]
    This means that the symmetry transformation \(g\) of \(G\) must be implemented by a unitary or antiunitary operator \(U_g\) on the Hilbert space \(\mathcal{H}\) of the system such that
    \[
        \ket{^g \Psi} = U_g \ket{\Psi}, \quad ^g O = U_g O U_g^{-1}.
    \]
\end{example}

In this example we have highlighted the importance of the \textbf{Wigner's theorem}:
\begin{theorem}[Wigner's theorem]
    Let \(\mathcal{H}\) be a Hilbert space and let \(G\) be a symmetry group of a quantum system represented on \(\mathcal{H}\). Then, for each symmetry transformation \(g\) of \(G\), there exists a unitary or antiunitary operator \(U_g\) on \(\mathcal{H}\) such that
    \begin{equation}
        \begin{dcases}
            ^g \Psi = U(g) \Psi,        \\
            ^g O = U_g O U_g^{\dagger}, \\
        \end{dcases}
        \label{eq:unitary_representation_symmetry_group_definition}
    \end{equation}
    since \(U_g\) is unitary or antiunitary, \(U_g^{-1} = U_g^{\dagger}\), and so we can write the transformation of observables as \(O \mapsto U_g O U_g^{\dagger}\). This implies that the expectation value of any observable \(O\) in any state \(\Psi\) is invariant under the action of the symmetry group \(G\).
\end{theorem}
This is a fundamental result in quantum theory since it allows us to represent symmetry transformations as unitary or antiunitary operators on the Hilbert space of the system.

Recall that a unitary operator \(U\) on a Hilbert space \(\mathcal{H}\) is a linear operator that preserves the Hilbert inner product,
\[
    \begin{aligned}
        U (\Phi\,a ,\, \Phi'\,a')  & = U (\Phi) a,\, U(\Phi')a^{\prime} \\
        \bra{U \Phi} \ket{U \Phi'} & = \bra{\Phi} \ket{\Phi'},
    \end{aligned}
\]
while an antiunitary operator \(A\) on a Hilbert space \(\mathcal{H}\) is an antilinear operator that preserves the Hilbert inner product,
\[
    \begin{aligned}
        A (\Phi\,a ,\, \Phi'\,a')  & = A (\Phi) a^*,\, A(\Phi')a^{\prime *} \\
        \bra{A \Phi} \ket{A \Phi'} & = \bra{\Phi} \ket{\Phi'}^*.
    \end{aligned}
\]
in both cases the adjoint operator \(U^{\dagger}\) (or \(A^{\dagger}\)) of a unitary (or antiunitary) operator \(U\) (or \(A\)) are equal to the inverse operator \(U^{-1}\) (or \(A^{-1}\)):
\[
    \begin{dcases}
        \bra{\Psi^{\prime}}\ket{U \Psi} = \bra{U^{\dagger} \Psi^{\prime}} \ket{\Psi}, \\
        \bra{\Psi^{\prime}}\ket{A \Psi} = (\bra{A^{\dagger} \Psi^{\prime}} \ket{\Psi})^*,
    \end{dcases} \quad \implies \quad \begin{dcases}
        U^{-1} = U^{\dagger}, \\
        A^{-1} = A^{\dagger},
    \end{dcases}
\]
respectively for the linear and antilinear case. In what follows, we leave aside the antiunitary case, which occurs only for symmetry transformations involving time reversal, and concentrate on the unitary one. Letting the symmetry transformation \(g\) vary in the group \(G\), we obtain a family of unitary operators \(U(g)\) contained in the operator algebra \(L(\mathcal{H})\).

Let us come back to relation \eqref{eq:unitary_representation_symmetry_group_definition}\(_1\). We recall again that state vectors are normalized vectors defined up to a multiplicative phase factor. The content of \eqref{eq:unitary_representation_symmetry_group_definition}\(_1\) is that phase choice can be adjusted so that for any symmetry transformation \(g\) the state correspondence \(\Psi \mapsto ^g \Psi\) is codified by a unitary operator \(U(g)\). This however does not completely fix \(U(g)\): we are still free to replace \(U(g)\) by
\begin{equation}
    \hat{U}(g) = \alpha(g) U(g),
    \label{eq:unitary_representation_phase_redefinition}
\end{equation}
where \(\alpha(g)\) is a complex number of unit modulus, which is called a \textbf{phase factor}. This involves a phase redefinition
\[
    ^g \Psi \to \alpha(g) \, ^g \Psi,
\]
in \eqref{eq:unitary_representation_symmetry_group_definition}\(_1\), which does not change the physical content of the transformation since states are defined up to a multiplicative phase factor. By a suitable choice of the phase factors \(\alpha(g)\) many relations simplyfy.

It is natural and possible to require
\[
    ^e \Psi = \Psi, \quad \forall \Psi \in \mathcal{H},
\]
so that the neutral element \(e\) of the symmetry group \(G\) is represented by the identity operator
\[
    U(e) = \mathbb{I},
\]
on the Hilbert space \(\mathcal{H}\) of the system. With this choice, we can now analyze the implications of the group properties of \(G\) on the family of unitary operators \(U(g)\). Considering two symmetry transformations \(g\), \(g'\) of \(G\), and a physical state \(\Psi\), we have \(^{g^{\prime}}(^g \Psi)\) and \(^{g^{\prime} \cdot g} \Psi\) defined only up to a conventional choice of a multiplicative \textbf{phase factor} as seen above, one cannot expect that a relation of the form \(^{g^{\prime}}(^g \Psi) = ^{g^{\prime} \cdot g} \Psi\) to hold strictly in general, but only up to some phase factor
\[
    \gamma(g, g^{\prime}, \Psi) \in \mathbb{C}, \quad |\gamma(g, g^{\prime}, \Psi)| = 1,
\]
depending on the symmetry transformations \(g\), \(g'\) and the state \(\Psi\). Thus we can write
\[
    ^{g^{\prime}}(^g \Psi) = ^{g^{\prime} \cdot g} \Psi \, \gamma(g, g^{\prime}, \Psi).
\]
Let us then express an important result concerning the phase factor \(\gamma(g, g^{\prime}, \Psi)\) appearing in the above relation.
\begin{proposition}
    The phase factor \(\gamma(g, g^{\prime}, \Psi)\) can be chosen to be independent of the state \(\Psi\), thus obtaining a relation of the form
    \begin{equation}
        U(g \cdot g^{\prime}) = \gamma(g, g^{\prime}) U(g) U(g^{\prime}),
        \label{eq:projective_representation_symmetry_group_definition}
    \end{equation}
    which is called a \textbf{projective representation} of the symmetry group \(G\). If the phase factor \(\gamma(g, g^{\prime})\) can be chosen to be 1 for all \(g\), \(g'\) in \(G\), then we have a \textbf{unitary representation} of the symmetry group \(G\).
\end{proposition}
\begin{proof}
    Proof...\TODO{Complete the proof if needed}
\end{proof}

The phase factors \(\gamma(g, g^{\prime})\) appearing in the projective representation of the symmetry group \(G\) are not arbitrary, but must satisfy some consistency conditions, for example the \textit{2-cocyclicity}:

\begin{proposition}
    Given three symmetry transformations \(g\), \(g'\), \(g''\) of \(G\), the following relation must hold:
    \begin{equation}
        \gamma(g, g' \cdot g'') \gamma(g', g'') = \gamma(g \cdot g', g'') \gamma(g, g'),
        \label{eq:2-cocyclicity_definition}
    \end{equation}
    where functions \(\gamma : G \times G \to U(1)\) satisfying the 2-cocyclicity condition are called \textbf{2-cocycles} of \(G\).
\end{proposition}
The 2-cocycle condition is a consequence of the associativity property of the group multiplication of \(G\).
\begin{proof}
    Proof...\TODO{Complete the proof if needed}
\end{proof}

\begin{proposition}
    The \(G\) 2-cocycle \(\gamma(g, g')\) enjoys a few special properties, which albeit technical are often useful at various points:
    \begin{equation}
        \begin{aligned}
            \gamma(e, g)      & = \gamma(g, e) = 1,  \\
            \gamma(g, g^{-1}) & = \gamma(g^{-1}, g),
        \end{aligned}
        \label{eq:2-cocycles_properties}
    \end{equation}
    for any symmetry transformation \(g\) of \(G\).
\end{proposition}
The first property is a consequence of the neutrality property of the group multiplication of \(G\), while the second one is a consequence of the invertibility property of the group multiplication of \(G\).
\begin{proof}
    Proof...\TODO{Complete the proof if needed}
\end{proof}

Relation \eqref{eq:2-cocyclicity_definition} entails also that, for any symmetry transformation \(g\), the operators \(U(g^{-1})\) and \(U(g)^{-1}\) differ by a \(g\)-dependent phase factor. Specifically, the relation between them takes the form
\[
    U(g^{-1}) = \varpi(g) U(g)^{-1}, \quad \varpi(g) = \gamma(g, g^{-1})^{-1}.
\]
\begin{proof}
    Proof...\TODO{Complete the proof if needed}
\end{proof}

In general it is not possible to absorb the phase \(\gamma(g^{\prime},g)\) by redefining the operators \(U(g)\) into the operators \(\hat{U}(g)\) given in \eqref{eq:unitary_representation_phase_redefinition} with \(\alpha(g)\) a suitable phase function. By doing so \(\gamma(g^{\prime},g)\) gets replaced by
\begin{equation}
    \hat{\gamma}(g, g') = \gamma(g, g') \alpha(g \cdot g') \alpha(g)^{-1} \alpha(g')^{-1},
    \label{eq:2-cocycle_phase_factor_definition}
\end{equation}
and in general this cannot be made equal to 1 by a judicious choice of \(\alpha(g)\). A phase function \(\gamma(g^{\prime},g)\) of the form
\begin{equation}
    \gamma(g, g') = \alpha(g \cdot g') \alpha(g)^{-1} \alpha(g')^{-1},
    \label{eq:2-coboundary_definition}
\end{equation}
for some function \(\alpha : G \to U(1)\) is called a \textbf{2-coboundary} of \(G\).

\begin{proposition}
    A G 2-coboundary is a G 2-cocycle, that is, any function \(\gamma(g, g')\) of the form \eqref{eq:2-coboundary_definition} satisfies the 2-cocyclicity condition \eqref{eq:2-cocyclicity_definition}.
\end{proposition}
\begin{proof}
    Proof...\TODO{Complete the proof if needed}
\end{proof}

The G 2–cocycle \(\hat{\gamma}(g^{\prime},g)\) in \eqref{eq:2-cocycle_phase_factor_definition} can be rendered equal to 1 precisely when the G 2–cocycle \(\gamma(g^{\prime},g)\) is a G 2–coboundary of the form \eqref{eq:2-coboundary_definition}. This may or may not happen. The study of these matters is the object of a branch of algebra known as group cohomology. Some special results will be obtained below. At any rate, if \(\hat{\gamma}(g^{\prime},g)\) can be made equal to 1, the representation \(\hat{U}(g)\) becomes
genuinely unitary, so that
\[
    \hat{U} (g \cdot g') = \hat{U}(g) \hat{U}(g'),
\]
for all symmetry transformations \(g\), \(g'\) of \(G\). Then the unitary family \(U(g)\) forms a genuine \textbf{unitary representation} of the symmetry group \(G\).

\section{Dynamical Symmetry Groups}

Suppose that for a symmetry group \(G\), the following is observed: by acting with any symmetry transformations of \(G\) on the devices that prepare the states of the system but not the apparatus that measures energy, the results of energy measurements do not change. The symmetry is then called \textbf{dynamical}, because \textit{energy is the infinitesimal generator of time shifts}.

We note that also by acting with any symmetry transformations of \(G\) on both the devices preparing the states of the system and the apparatus measuring energy, the results of energy measurements do not change. Therefore, when \(G\) is a dynamical symmetry group, then the system’s Hamiltonian \(H\) satisfies for any symmetry element \(g\) of \(G\) the relation
\begin{equation}
    ^g H = H.
    \label{eq:dynamical_symmetry_hamiltonian_invariance_definition}
\end{equation}
If \(G\) acts on th Hilbert space through a unitary representation \(U(g)\), then the invariance of the Hamiltonian under the action of \(G\) can be written as
\begin{equation}
    U(g) H U(g)^{\dagger} = H, \quad \forall g \in G,
    \label{eq:dynamical_symmetry_hamiltonian_invariance_unitary_representation_definition}
\end{equation}
by virtue of Wigner's theorem \eqref{eq:unitary_representation_symmetry_group_definition}\(_2\). This means that the Hamiltonian \(H\) commutes with all the unitary operators \(U(g)\) representing the symmetry transformations of \(G\):
\[
    \left[ H,\, U(g) \right] = 0, \quad \forall g \in G.
\]
This is a non trivial property severely constraining the form of \(H\) and thus providing useful structural information about the form of the interactions the system is involved in. Let us illustrate this point with a few examples taken from quantum theory.

\begin{example}[The Hydrogen Atom]
    The hydrogen atom is a system constituted by a proton and an electron interacting through the Coulomb potential. The Hamiltonian of the hydrogen atom is given by
    \[
        H = \frac{\mathbf{p}^2}{2m} V(\mathbf{q}).
    \]
    Direct observation shows that all proper rotations of the devices preparing the states of the atom leave the results of all energy measurements unchanged. Therefore, the proper rotation group \(\mathrm{SO}(3)\) is dynamical. The atom has spherical symmetry, so from \eqref{eq:dynamical_symmetry_hamiltonian_invariance_unitary_representation_definition} we have that the Hamiltonian ha to remain unchanged under the action of any proper rotation \(\mathbf{R}\) of \(\mathrm{SO}(3)\):
    \[
        U(\mathbf{R}) H U(\mathbf{R})^{\dagger} = H, \quad \forall \mathbf{R} \in \mathrm{SO}(3).
    \]
    It is sensible to assum that \(^{\mathbf{R}}\mathbf{q} = \mathbf{R}^{-1}\mathbf{q}\) and \(^{\mathbf{R}}\mathbf{p} = \mathbf{R}^{-1}\mathbf{p}\) under the action of a proper rotation. Then, we can write
    \[
        U(\mathbf{R}) \mathbf{q} U(\mathbf{R})^{\dagger} = \mathbf{R}^{-1}\mathbf{q}, \quad U(\mathbf{R}) \mathbf{p} U(\mathbf{R})^{\dagger} = \mathbf{R}^{-1}\mathbf{p},
    \]
    which means that the position and momentum operators transform as vectors under the action of \(\mathrm{SO}(3)\).\footnote{This is a consequence of the fact that they transform according to the fundamental representation of \(\mathrm{SO}(3)\): \(\mathbf{q} \mapsto \mathbf{R}^{-1}\mathbf{q}\), \(\mathbf{p} \mapsto \mathbf{R}^{-1}\mathbf{p}\).}

    Using this last relation along with the invariance of the Hamiltonian under the action of \(\mathrm{SO}(3)\) implies that
    \[
        \frac{(\mathbf{R}^{-1}\mathbf{p})^2}{2m} + V(\mathbf{R}^{-1}\mathbf{q}) = \frac{\mathbf{p}^2}{2m} + V(\mathbf{q}), \quad \forall \mathbf{R} \in \mathrm{SO}(3).
    \]
    In particular, knowing that the rotations preserve the length of vectors, we have \((\mathbf{R}^{-1}\mathbf{p})^2 = \mathbf{p}^2\), so that the invariance of the Hamiltonian under the action of \(\mathrm{SO}(3)\) reduces to
    \[
        V(\mathbf{R}^{-1} \mathbf{q}) = V(\mathbf{q}) \quad \implies \quad V(\mathbf{q}) = v(|\mathbf{q}|),
    \]
    so that the potential \(V\) must be a function of the length of the position vector \(v(|\mathbf{q}|)\) only. This is a non trivial constraint on the form of the potential \(V\) and thus on the form of the interactions the system is involved in; it is not sufficient to yield the Coulombic expression of the potential, but it is a necessary condition for it. The spherical symmetry of the hydrogen atom is a consequence of the dynamical symmetry of \(\mathrm{SO}(3)\).
\end{example}

\begin{example}[A System of Identical Particles]\TODO{Complete the example}

\end{example}

\begin{example}[Conduction Electrons in a Crystal Lattice]\TODO{Complete the example}

\end{example}

\subsection{Energy Eigenvalues Classification}

The energy spectrum of a quantum system typically contains many energy values. The problem then arises of devising a classification method for the energy eigenvalues. There does not exist any a priori ideal classification scheme. However, symmetry provides a natural one. In this section, we show that when the system is characterized by a dynamical symmetry group \(G\), the energy eigenvalues can be efficiently classified based on the unitary representations of \(G\) according to which the energy eigenstates transform under the action of \(G\).

We assume that the system’s Hamiltonian \(H\) has a \textbf{completely discrete spectrum} and that each energy eigenvalue has \textbf{finite degeneracy}. Let \(E\) be an energy eigenvalue of \(H\) and let \(\mathcal{H}_E\) be the corresponding eigenspace, which is \(d_E\)-dimensional given a degeneracy \(d_E\); this subspace is spanned by all vectors \(\Psi\) which respect the eigenvalue equation
\[
    H \Psi = E \Psi.
\]
Let now \(\Psi\) be a vector of \(\mathcal{H}_E\), then for any symmetry transformation \(g\) of \(G\) we have
\[
    H U(g)\Psi = U(g) H \Psi = E U(g) \Psi,
\]
as a direct consequence of \eqref{eq:dynamical_symmetry_hamiltonian_invariance_unitary_representation_definition}. This means that \(U(g)\Psi\) is also an eigenvector of \(H\) with the same eigenvalue \(E\). Thus, the eigenspace \(\mathcal{H}_E\) is invariant under the action of \(G\). The unitary operators \(U(g)\) representing the symmetry transformations of \(G\) on the Hilbert space \(\mathcal{H}\) of the system restrict to unitary operators on the eigenspace \(\mathcal{H}_E\), thus providing a unitary representation of \(G\) on \(\mathcal{H}_E\). The energy eigenvalue \(E\) can be classified according to the unitary representation of \(G\) according to which the energy eigenstates in \(\mathcal{H}_E\) transform under the action of \(G\).

\subsubsection{Unitary Representations on the Eigenspaces}

If we now consider \( \{ \Phi_i \} \) an ON basis for the subspace \(\mathcal{H}_E\), then these vectors respect the eigenvalue equation and the orthonormality condition
\[
    \begin{dcases}
        H \Phi_i = E \Phi_i, \\
        \bra{\Phi_i} \ket{\Phi_j} = \delta_{ij}.
    \end{dcases}
\]
Since the unitary representation of \(G\) on \(\mathcal{H}_E\) is given by the restriction of the unitary representation of \(G\) on \(\mathcal{H}\), we have that for any symmetry transformation \(g\) of \(G\) and any vector \(\Phi_i\) of the ON basis \(\{\Phi_i\}\) of \(\mathcal{H}_E\), the transformed vector \(U(g)\Phi_i\) is a linear combination of the vectors of the ON basis \(\{\Phi_i\}\):
\begin{equation}
    U(g)\Phi_i = \sum_{j=1}^{d_E} \Phi_j D_{ji}(g),
    \label{eq:unitary_representation_eigenspace_linear_combination_definition}
\end{equation}
where \(D_{ji}(g)\) are the complex, \(g\)-dependent matrix elements of the unitary representation of \(G\) on \(\mathcal{H}_E\). We aim to prove that the family of matrices \(D(g)\) with elements \(D_{ji}(g)\) forms a unitary representation of the symmetry group \(G\) on the eigenspace \(\mathcal{H}_E\).
\begin{proposition}
    For each element \(g\) of the symmetry group \(G\), there exists a \(d_E \times d_E\) complex matrix \(D(g) = (D_{ji}(g))\) which is unitary:
    \begin{equation}
        \begin{aligned}
            D(g \cdot g^{\prime}) & = D(g) D(g^{\prime}),               \\
            D(g)^{\dagger} D(g)   & = D(g) D(g)^{\dagger} = \mathbb{I},
        \end{aligned}
        \label{eq:unitary_representation_eigenspace_definition}
    \end{equation}
    for all \(g\), \(g'\) in \(G\). The family of unitary matrices \(D(g)\) forms a unitary representation of the symmetry group \(G\) on the eigenspace \(\mathcal{H}_E\) corresponding to the energy eigenvalue \(E\).
\end{proposition}

\begin{proof}
    As \(U(g)\) is a unitary representation of \(G\) on the Hilbert space \(\mathcal{H}\) of the system, we know that for any \(g\), \(g'\) in \(G\) we have
    \[
        U(g \cdot g^{\prime}) \Phi_i = U(g) U(g^{\prime}) \Phi_i.
    \]
    Since the vectors \(\Phi_i\) form a basis for the eigenspace \(\mathcal{H}_E\), the above equation implies that the matrix elements of \(U(g)\) in this basis satisfy the same multiplication rule as the matrices \(D(g)\), from \eqref{eq:unitary_representation_eigenspace_linear_combination_definition}:
    \[
        U(g \cdot g^{\prime}) \Phi_i = \sum_{j=1}^{d_E} \Phi_j D_{ji}(g \cdot g^{\prime}),
    \]
    and
    \[
        \begin{aligned}
            U(g) U(g^{\prime}) \Phi_i & = U(g) \sum_{j} \Phi_j D_{ji}(g^{\prime}) = \sum_{j} U(g) \Phi_j D_{ji}(g^{\prime})                                                            \\
                                      & = \sum_{j} \left( \sum_{k} \Phi_k D_{kj}(g) \right) D_{ji}(g^{\prime}) = \sum_{k} \Phi_k \left( \sum_{j} D_{kj}(g) D_{ji}(g^{\prime}) \right).
        \end{aligned}
    \]
    Now if we put together the above two equations, we can prove the first relation in \eqref{eq:unitary_representation_eigenspace_definition}:
    \[
        D_{ji} (g \cdot g^{\prime}) = \sum_{k} D_{kj}(g) D_{ji}(g^{\prime}),
    \]
    which is the multiplication rule for the matrices \(D(g)\).

    The second relation in \eqref{eq:unitary_representation_eigenspace_definition} is a consequence of \eqref{eq:unitary_representation_eigenspace_linear_combination_definition} and the orthonormality of the basis \(\{\Phi_i\}\):
    \[
        \begin{aligned}
            \delta_{ij} & = \bra{\Phi_i} \ket{\Phi_j} = \bra{U(g) \Phi_i} \ket{U(g) \Phi_j}                                                               \\
                        & = \bra{\sum_{k} \Phi_k D_{ki}(g)} \ket{\sum_{l} \Phi_l D_{lj}(g)} = \sum_{k, l} D_{ki}(g)^* \bra{\Phi_k} \ket{\Phi_l} D_{lj}(g) \\
                        & = \sum_{k} D_{ik}(g)^{\dagger} D_{kj}(g) = (D(g)^{\dagger} D(g))_{ij},
        \end{aligned}
    \]
    where in the last steps we have transposed the indices and used the definition of the adjoint matrix \(D(g)^{\dagger}\). This proves that \(D(g)^{\dagger} D(g) = \mathbb{I}\).
\end{proof}

Having proven \eqref{eq:unitary_representation_eigenspace_definition}, we know \(g \mapsto D(g)\) to be a unitary representation of the symmetry group \(G\) on the eigenspace \(\mathcal{H}_E\) corresponding to the energy eigenvalue \(E\), and from \eqref{eq:unitary_representation_eigenspace_linear_combination_definition} we know the basis vectors \(\Phi_i\) to transform under the action of \(G\) through \(D(g)\).

The energy eigenvalue \(E\) can thus be classified according to the unitary representation \(D(g)\), where the energy eigenstates in \(\mathcal{H}_E\) transform under the action of \(G\). The problem is that there are infinitely many orthonormal bases of the energy eigenspace \(\mathcal{H}_E\) of \(E\). The representation \(D(g)\) depends on the chosen basis \(\{\Phi_i\}\). So \(D(g)\) by itself cannot be used to distinguish \(E\) from other energy eigenvalues. Let us analyze this point in greater detail.

\begin{proposition}
    Let \(\{\Phi_{\ i}^1\}\) and \(\{\Phi_{\ i}^2\}\) be two ON bases for the eigenspace \(\mathcal{H}_E\) corresponding to the energy eigenvalue \(E\). Then, let the corresponding unitary representations be \(D^1(g)\) and \(D^2(g)\), defined by \eqref{eq:unitary_representation_eigenspace_linear_combination_definition} with respect to the bases \(\{\Phi_{\ i}^1\}\) and \(\{\Phi_{\ i}^2\}\), respectively.

    There exist a \(d_E \times d_E\) unitary matrix \(A\) such that its elements \(A_{ij}\) relate the two bases \(\{\Phi_{\ i}^1\}\) and \(\{\Phi_{\ i}^2\}\) as
    \begin{equation}
        \Phi_{\ i}^2 = \sum_{j=1}^{d_E} \Phi_{\ j}^1 A_{ji},
        \label{eq:unitary_representation_eigenspace_basis_change_definition}
    \end{equation}
    and the two unitary representations \(D^1(g)\) and \(D^2(g)\) are related by the similarity transformation
    \begin{equation}
        D^2(g) = A^{-1} D^1(g) A,
        \label{eq:unitary_representation_eigenspace_similarity_transformation_definition}
    \end{equation}
    for all \(g\) in \(G\). Thus, the two unitary representations \(D^1(g)\) and \(D^2(g)\) are equivalent representations of the symmetry group \(G\).
\end{proposition}

\begin{proof}
    Using \eqref{eq:unitary_representation_eigenspace_linear_combination_definition}, we can develop \eqref{eq:unitary_representation_eigenspace_basis_change_definition} as
    \[
        \begin{aligned}
            U(g) \Phi_{\ i}^2 & = U(g) \sum_{j=1}^{d_E} \Phi_{\ j}^1 A_{ji} = \sum_{j=1}^{d_E} U(g) \Phi_{\ j}^1 A_{ji} = \sum_{j=1}^{d_E} \left( \sum_{k=1}^{d_E} \Phi_{\ k}^1 D_{kj}^1(g) \right) A_{ji} \\
                              & = \sum_{k=1}^{d_E} \Phi_{\ k}^1 \left( \sum_{j=1}^{d_E} D_{kj}^1(g) A_{ji} \right) = \sum_{k=1}^{d_E} \Phi_{\ k}^1 (D^1(g) A)_{ki},
        \end{aligned}
    \]
    on one hand, and
    \[
        \begin{aligned}
            U(g) \Phi_{\ i}^2 & = \sum_{k=1}^{d_E} \Phi_{\ k}^2 D_{ki}^2(g) = \sum_{k=1}^{d_E} \left( \sum_{j=1}^{d_E} \Phi_{\ j}^1 A_{jk} \right) D_{ki}^2(g)      \\
                              & = \sum_{j=1}^{d_E} \Phi_{\ j}^1 \left( \sum_{k=1}^{d_E} A_{jk} D_{ki}^2(g) \right) = \sum_{j=1}^{d_E} \Phi_{\ j}^1 (A D^2(g))_{ji}.
        \end{aligned}
    \]
    Putting together the above two equations, we have
    \[
        (D^1(g) A)_{ki} = (A D^2(g))_{ki}, \quad \text{i.e., } D^1(g) A = A D^2(g),
    \]
    which shows that the two unitary representations \(D^1(g)\) and \(D^2(g)\) are related by the similarity transformation \(D^2(g) = A^{-1} D^1(g) A\). This proves that the two unitary representations \(D^1(g)\) and \(D^2(g)\) are equivalent representations of the symmetry group \(G\).
\end{proof}

Thus two unitary representations \(D^1(g)\) and \(D^2(g)\) of the symmetry group \(G\) of the same dimension \(d_E\), related as in \eqref{eq:unitary_representation_eigenspace_similarity_transformation_definition}, for some unitary matrix \(A\), are said to be \textbf{unitarily equivalent}. Thus the unitary representation \(D(g)\) in \eqref{eq:unitary_representation_eigenspace_linear_combination_definition} is independent of the choice of the ON basis \(\{\Phi_i\}\) for the eigenspace \(\mathcal{H}_E\) up to unitary equivalence.

We can in this way characterize the energy eigenvalue \(E\) by a \textit{unitary equivalence class of unitary representations}\footnote{It is a class of equivalence of unitary representations, where the representations are unitarily equivalent.} of \(G\) and use this class to group \(E\) with other eigenvalues sharing the very same class. In this way, \textbf{a symmetry based classification of the energy spectrum is achieved}. Since however there are infinitely many representations of any given group, the classification scheme obtained may not be so useful. Fortunately, generic representations can be decomposed as combinations of a smaller set of more basic representations, which can then be used to refine the scheme.

\subsubsection{Reducibility and Irreducibility}

Suppose that we are able to chose the orthonormal basis \(\{\Phi_i\}\) of the energy eigenspace \(\mathcal{H}_E\) to be of indexed \(\Phi_{\alpha r}\) by two indices \(\alpha\), \(r\) in such a way that \eqref{eq:unitary_representation_eigenspace_linear_combination_definition} takes the form
\[
    U(g) \Phi_{\alpha r} = \sum_{s} \Phi_{\alpha s} D_{\alpha s r}(g),
\]
for every element \(g\) of \(G\), for certain \(g\)-dependent coefficients \(D_{\alpha s r}(g)\). Therefore, the vectors \(\Phi_{\alpha r}\) for fixed \(\alpha\) mix among themselves under the action of \(G\). This means that, for fixed \(g\), the matrix \(D(g)\) has the block diagonal form
\[
    D(g) = \begin{pmatrix}
        D_1(g) & 0      & \cdots & 0      \\
        0      & D_2(g) & \cdots & 0      \\
        \vdots & \vdots & \ddots & \vdots \\
        0      & 0      & \cdots & D_a(g)
    \end{pmatrix},
\]
and let us assume that the blocks \(D_{\alpha}(g)\) are of dimension \(d_{\alpha}\), which is the smallest non trivial dimension for each block (i.e. the blocks cannot be further decomposed into diagonal blocks by a furhter change of basis \(\Phi_{\alpha r}\)). Then, the unitary representation \(g \mapsto D(g)\) of \(G\) on \(\mathcal{H}_E\) is said to be \textbf{reducible}, and its blocks \(D_{\alpha}(g)\) are said to define an \textbf{irreducible} representations \(g \mapsto D_{\alpha}(g)\) of \(G\) (\textit{irreducible components}).

The energy eigenvalue \(E\) can then be classified through the irreducible representations \(D_{\alpha}(g)\) of \(G\). The former way of proceeding provides a classification scheme of the energy spectrum relying on a more restricted set of elements and so potentially more useful than that furnished by the latter.

We can resume the results obtained in this section as follows:
\begin{itemize}
    \item Each energy eigenvalue \(E\) is characterized by a set of irreducible unitary representations \(D_{\alpha}\) of \(G\), unique up to unitary equivalence, describing how the eigenvectors belonging to \(E\) transform under the action of \(G\).
    \item The energy eigenvalues can be grouped in subsets of eigenvalues sharing the same representation content \(D_{\alpha}\) providing a classification scheme of the energy spectrum.
\end{itemize}

Let us illustrate the above results with a few examples.

\begin{example}[The One-Dimensional Harmonic Oscillator]
    The parity transformation \(P\) is space reflection about the origin. Observation show that the parity group \(C_2 = \{\mathbb{I},P\}\) is a dynamical symmetry group of the one–dimensional harmonic oscillator.

    As well known, the energy eigenvalues of the one–dimensional harmonic oscillator are given by
    \[
        H \Phi_n = E_n \Phi_n. \quad E_n = \hbar \omega (n + 1/2),
    \]
    with \(n = 0, 1, 2, \ldots\), and are non degenerate. The energy eigenstates of the one–dimensional harmonic oscillator transform under the action of the parity transformation \(P\) as
    \[
        U(P) \Phi_n = \Phi_n (-1)^n.
    \]
    Thus the parity group \(C_2\) has two irreducible unitary representations: the trivial representation \(D_+(P) = 1\) and the sign representation \(D_-(P) = -1\). The energy eigenvalues of the one–dimensional harmonic oscillator are classified according to these two irreducible unitary representations of \(C_2\): the energy eigenvalues with even \(n\) transform according to the trivial representation \(D_+\), while those with odd \(n\) transform according to the sign representation \(D_-\). Thus we classify them as even or odd energy eigenvalues according to the parity of \(n\).
\end{example}

\begin{example}[The Hydrogen Atom]
    It can be shown that the proper rotation group \(\mathrm{SO}(3)\) is a dynamical symmetry group of the hydrogen atom. The energy eigenvalues of the hydrogen atom are known to be given by
    \[
        H \Phi_{n l m} = E_n \Phi_{n l m}, \quad E_n = -\frac{m c^2 \alpha^2}{2 n^2}.
    \]
    The energy eigenvalues of the hydrogen atom are degenerate, with degeneracy \(d_{E_n} = n^2\):
    \begin{enumerate}
        \item \(n = 0, 1, 2, \ldots\),
        \item \(l = 0, 1, \ldots, n-1\),
        \item \(m = -l, -l+1, \ldots, l-1, l\),
    \end{enumerate}
\end{example}
and the eigenstates \(\Phi_{n l m}\) are also eigenstates of the total angular momentum operator \(L^2\) and of the component along the quantization axis of the angular momentum operator \(L_z\):
\[
    \begin{dcases}
        L^2 \Phi_{n l m} = \hbar^2 l(l+1) \Phi_{n l m}, \\
        L_0 \Phi_{n l m} = \hbar m \Phi_{n l m}.
    \end{dcases}
\]
Furthermore, the other spherical coordinates of the angular momentum \(L_{\pm}\) act on the eigenstates \(\Phi_{n l m}\) as
\[
    L_{\pm} \Phi_{n l m} = \hbar \sqrt{l(l+1) - m(m \pm 1)} \Phi_{n l (m \pm 1)},
\]
thus acting as ladder operators on the index \(m\).

Since \(\mathbf{L}\) is the \textbf{infinitesimal generator} of the proper rotation group \(\mathrm{SO}(3)\) (which will be discussed in the next section \ref{sec:lie}), the energy eigenstates \(\Phi_{n l m}\) transform under the action of \(\mathrm{SO}(3)\) according to the irreducible unitary representation \(D_l(\mathbf{R})\) of \(\mathrm{SO}(3)\), of dimension \(d_l = 2l + 1\):
\[
    U(\mathbf{R}) \Phi_{n l m} = \sum_{m'=-l}^{l} \Phi_{n l m'} D_{l m' m}(\mathbf{R}),
\]
so that \(D_l(\mathbf{R})\) is \((2l+1)\times(2l+1)\) dimensional, and the mappings \(\mathbf{R} \mapsto D_l(\mathbf{R})\) can be shown to be unitary irreducible representations of \(\mathrm{SO}(3)\). Thus for each energy eigenvalue \(E_n\) of the hydrogen atom, the energy eigenstates \(\Phi_{n l m}\) transform under the action of \(\mathrm{SO}(3)\) according to the \(n\) irreducible \((2l+1)\)-dimensional unitary representation \(D_{l=0,1,\ldots,n-1}\) of \(\mathrm{SO}(3)\). In this way, the symmetry based classification of the energy spectrum is achieved for the hydrogen atom.

\begin{example}[The Sulfur Trioxide Molecule SO\(_3\)]\TODO{put image of the molecule}
    The molecule of sulfur trioxide SO\(_3\) has the shape of a planar equilateral triangle, with the sulfur atom at the center and the three oxygen atoms at the vertices. The symmetry group of this molecule is the dihedral group \(D_3\), which is a subgroup of \(\mathrm{SO}(3)\).

    A planar equilateral triangle has the following symmetries:
    \begin{itemize}
        \item Three rotations of \(120^{\circ}\) about the axis perpendicular to the plane of the triangle, which we call \(C_3\) rotations:
              \[
                  0^{\circ} : \begin{pmatrix}
                      123        \\
                      \downarrow \\
                      123
                  \end{pmatrix}, \quad 120^{\circ} : \begin{pmatrix}
                      123        \\
                      \downarrow \\
                      231
                  \end{pmatrix}, \quad 240^{\circ} : \begin{pmatrix}
                      123        \\
                      \downarrow \\
                      312
                  \end{pmatrix}.
              \]
        \item Three reflections about the axes of symmetry of the triangle, which we call \(\sigma_v\) reflections:
              \[
                  \sigma_v^1 : \begin{pmatrix}
                      123        \\
                      \downarrow \\
                      132
                  \end{pmatrix}, \quad \sigma_v^2 : \begin{pmatrix}
                      123        \\
                      \downarrow \\
                      321
                  \end{pmatrix}, \quad \sigma_v^3 : \begin{pmatrix}
                      123        \\
                      \downarrow \\
                      213
                  \end{pmatrix}.
              \]
    \end{itemize}
    Putting together the above symmetries, we find a total of six symmetries, which form the dihedral group \(D_3\), which coincides with the symmetric group \(S_3\) of permutations of three objects. The energy eigenvalues of the sulfur trioxide molecule SO\(_3\) can be classified according to the irreducible unitary representations of \(D_3\).

    This group has three \textit{unitary equivalence classes of irreducible unitary representations}: the trivial representation \(D_+\), the sign representation \(D_-\) and a two-dimensional representation \(D_2\). They are of dimensions \(d_+ = 1\), \(d_- = 1\) and \(d_2 = 2\), respectively, and we can deduce their dimensionality through the following theorem:
    \begin{theorem}
        The sum of the squares of the dimensions \(d_{\alpha}\) of the irreducible unitary representations \(D_{\alpha}\) of a finite non-abelian group \(G\) is equal to the order \(|G|\) (the number of elements) of the group:
        \[
            \sum_{\alpha} d_{\alpha}^2 = |G|.
        \]
    \end{theorem}
    In our case the only way to satisfy the above relation is to have \(d_+^2 + d_-^2 + d_2^2 = 1^2 + 1^2 + 2^2 = 6\), which is the order of \(D_3\). The energy eigenvalues of the sulfur trioxide molecule SO\(_3\) can be classified according to these three irreducible unitary representations of \(D_3\).
\end{example}
